\section{Deletion}
\label{sec:deletion}

Deletion follows the same three stages as insertion. The key is removed the way
an ordinary binary search tree removes it, the procedure walks back toward the
root updating heights, and unbalanced nodes are repaired by rotation. The
difference lies in the third stage, and it is the reason deletion deserves its
own section: the walk cannot stop at the first repair.

\begin{algorithm}[t]
\caption{Deletion. A node with two children is removed by copying its in-order
successor into it and deleting the successor instead, which always has at most
one child.}
\label{alg:delete}
\SetKwFunction{Del}{Delete}
\SetKwFunction{Reb}{Rebalance}
\SetKwFunction{Min}{Minimum}
\Fn{\Del{$v$, $k$}}{
  \lIf{$v = \textbf{nil}$}{\KwRet \textbf{nil}}
  \uIf{$k < v.\mathit{key}$}{
    $v.\mathit{left} \gets$ \Del{$v.\mathit{left}$, $k$}\;
  }
  \uElseIf{$k > v.\mathit{key}$}{
    $v.\mathit{right} \gets$ \Del{$v.\mathit{right}$, $k$}\;
  }
  \Else{
    \lIf{$v.\mathit{left} = \textbf{nil}$}{\KwRet $v.\mathit{right}$}
    \lIf{$v.\mathit{right} = \textbf{nil}$}{\KwRet $v.\mathit{left}$}
    $s \gets$ \Min{$v.\mathit{right}$} \Comment*[r]{in-order successor}
    $v.\mathit{key} \gets s.\mathit{key}$\;
    $v.\mathit{right} \gets$ \Del{$v.\mathit{right}$, $s.\mathit{key}$}\;
  }
  \KwRet \Reb{$v$}\;
}
\end{algorithm}

\subsection{Why the repair propagates}

\Cref{thm:insertone} rests on one fact: after an insertion is repaired, the
subtree has the height it started with, so the parent notices nothing. Deletion
breaks that fact.

\begin{lemma}[The repaired subtree may shrink]
\label{lem:shrink}
Let $z$ be an unbalanced node after a deletion, with the subtree at $z$ having
height $h$ before the deletion. After the prescribed repair, the subtree rooted
at the new position has height $h$ or $h-1$, and the case $h-1$ can occur.
\end{lemma}

\begin{proof}
Suppose the right subtree of $z$ lost a level, so $\bfac(z) = +2$, and write $y$
for the left child of $z$. Before the deletion the subtree at $z$ had height $h$,
so $y$ has height $h-1$ and the right subtree of $z$ now has height $h-3$.

Consider the \case{ll} case, in which $\bfac(y) = +1$. Then $L(y)$ has height
$h-2$ and $R(y)$ has height $h-3$. The right rotation at $z$ puts $y$ on top with
$L(y)$ of height $h-2$ on the left, and on the right the node $z$ carrying $R(y)$
and the shortened right subtree, both of height $h-3$, hence of height $h-2$. The
rotated subtree therefore has height $h-1$, one less than before.

The case $\bfac(y) = 0$ gives the other outcome. Then $L(y)$ and $R(y)$ both have
height $h-2$, the rotation leaves $z$ with a subtree of height $h-2$ on each
side, and the result has height $h$ once more. Both outcomes are therefore
reachable, and the mirror cases behave identically.
\end{proof}

When the repaired subtree comes out one level shorter, the parent of $z$ has lost
a level on one side without gaining one on the other. Its own balance factor may
now reach $\pm 2$, and the same may happen to its parent, and so on. A single
deletion can therefore set off repairs at every level between the deleted node
and the root.

\subsection{A cascade in the smallest tree that permits one}

The trees in which this happens are exactly the minimum-size trees of
\Cref{sec:height}. They have no slack anywhere. Every node in the minimum-size
tree of height $h$ has a balance factor of exactly $\pm 1$, so removing any node
at all pushes some ancestor out of range.

\Cref{fig:cascade} shows the minimum-size tree of height $7$, which holds $54$
nodes by \Cref{tab:nh}. Removing the marked leaf drops the count to $53$. By
\Cref{eq:nrec} a tree of $53$ nodes cannot have height $7$, since $\Nmin(7) = 54$,
so the height must fall, and the repairs that make it fall run at three separate
levels. Each rotation shortens its subtree by one level in the manner of
\Cref{lem:shrink}, which is precisely what exposes the next ancestor.

\begin{figure}[t]
  \centering
  % The deletion cascade, on the smallest tree that can have one at this height.
%
% The drawing macro names its nodes f1, f2, ... in pre-order, so the marked
% positions below are not decoration placed by eye: f34 is the leaf that is
% removed, and f31, f23 and f2 are the three ancestors that rebalance, in that
% order. Every one of them is a left-left case.
%
% There is no slack anywhere in this tree. Every node has a balance factor of
% exactly one, so losing a single leaf is felt all the way up.
\begin{tikzpicture}[x = 1mm, y = 1mm]

  \setavlscale{6.6}{8.4}
  \resetavl
  \drawfib{7}{0}{0}

  % the leaf that goes
  \node[circle, draw = rkink, line width = 0.9pt, inner sep = 0pt,
        minimum size = 3.4mm] (victim) at (f34) {};
  \node[tlabel, below = 2mm of victim, anchor = north] {removed};

  % the three ancestors that rebalance, in the order they fire
  \foreach \nd/\num in {f31/1, f23/2, f2/3} {
    \node[circle, draw = rkink, fill = rwash, line width = 0.9pt,
          inner sep = 0pt, minimum size = 4.2mm] (r\num) at (\nd) {};
    \node[font = \scriptsize, above right = -0.6mm and 0.8mm of r\num] {\num};
  }

  % the height bracket on the left
  \draw[decorate, decoration = {brace, amplitude = 4pt}, draw = rmid]
    (-72, -58.8) -- (-72, 0)
    node[midway, left = 6pt, tlabel, rotate = 90, anchor = south]
      {height 7, before};

  \node[tannot, align = center] at (0, -70)
    {54 nodes, and every balance factor equal to one.\\
     Removing the marked leaf rebalances three ancestors and leaves height 6.};

\end{tikzpicture}

  \caption{A deletion cascade in the minimum-size tree of height $7$. Removing
  one leaf forces rebalancing at the three ringed ancestors, in the order shown,
  and the height of the tree falls from $7$ to $6$.}
  \label{fig:cascade}
\end{figure}

This example also shows why the cascade cannot be dismissed as a pathological
construction. The minimum-size trees are the extremal case for the height bound
of \Cref{thm:heightbound}, so any tree that comes close to the bound is close to
one of them, and is therefore close to having no slack.

\subsection{Cost}

The cascade is bounded, because the walk visits each level at most once.

\begin{theorem}[Deletion cost]
\label{thm:deletebound}
A deletion from an AVL tree of $n$ nodes performs at most $h$ rebalancing steps
and at most $2h$ single rotations, where $h$ is the height of the tree, and runs
in $O(\log n)$ time.
\end{theorem}

\begin{proof}
The removal descends one level at a time to the node holding the key, and if that
node has two children, the procedure descends further to its in-order successor.
Both descents visit at most $h$ nodes. The upward walk retraces the same path, so
it examines at most $h$ ancestors, and each ancestor is rebalanced at most once.
By \Cref{tab:cases} a rebalancing step performs one single rotation in the
\case{ll} and \case{rr} cases and two in the \case{lr} and \case{rl} cases,
giving at most $2h$ single rotations in total.

For the running time, the descent and the walk each cost $O(h)$, and by
\Cref{lem:rotcost} each rotation costs $O(1)$, so the rotations cost $O(h)$
together. The total is $O(h)$, which is $O(\log n)$ by \Cref{thm:heightbound}.
\end{proof}

\Cref{tab:insdel} sets the two update operations side by side. Both are
logarithmic, and they differ in the constant work they do on the way back up:
insertion stops at the first repair, deletion may repair at every level.

\begin{table}[t]
  \centering
  \caption{Rebalancing work for the two update operations, for a tree of height
  $h$. The asymptotic cost is the same; the number of repairs is not.}
  \label{tab:insdel}
  \begin{tabular}{@{}lccc@{}}
    \toprule
    Operation & Rebalancing steps & Single rotations & Total time \\
    \midrule
    Insertion & at most $1$  & at most $2$  & $O(\log n)$ \\
    Deletion  & at most $h$  & at most $2h$ & $O(\log n)$ \\
    \bottomrule
  \end{tabular}
\end{table}

The bound of \Cref{thm:deletebound} is a worst case, and \Cref{sec:eval} shows
that the average behaviour is far from it.
